\documentclass[10pt,a4paper]{amsart}
\usepackage{amsmath,amssymb,amsthm,graphicx}
\usepackage{verbatim}
\newenvironment{reference}{\comment}{\endcomment}
%\newenvironment{reference}{}{}
\newenvironment{slogan}{\comment}{\endcomment}
\newenvironment{history}{\comment}{\endcomment}

% For commutative diagrams we use Xy-pic
\usepackage[all]{xy}

% We use 2cell for 2-commutative diagrams.
\xyoption{2cell}
% Theorem environments.
%
\theoremstyle{plain}
\newtheorem{theorem}[subsection]{Theorem}
\newtheorem{proposition}[subsection]{Proposition}
\newtheorem{lemma}[subsection]{Lemma}

\theoremstyle{definition}
\newtheorem{definition}[subsection]{Definition}
\newtheorem{example}[subsection]{Example}
\newtheorem{exercise}[subsection]{Exercise}
\newtheorem{situation}[subsection]{Situation}

\theoremstyle{remark}
\newtheorem{remark}[subsection]{Remark}
\newtheorem{remarks}[subsection]{Remarks}

\numberwithin{equation}{subsection}

% Macros
%
\def\lim{\mathop{\mathrm{lim}}\nolimits}
\def\colim{\mathop{\mathrm{colim}}\nolimits}
\def\Spec{\mathop{\mathrm{Spec}}}
\def\Hom{\mathop{\mathrm{Hom}}\nolimits}
\def\Ext{\mathop{\mathrm{Ext}}\nolimits}
\def\SheafHom{\mathop{\mathcal{H}\!\mathit{om}}\nolimits}
\def\SheafExt{\mathop{\mathcal{E}\!\mathit{xt}}\nolimits}
\def\Sch{\mathit{Sch}}
\def\Mor{\mathop{\mathrm{Mor}}\nolimits}
\def\Ob{\mathop{\mathrm{Ob}}\nolimits}
\def\Sh{\mathop{\mathit{Sh}}\nolimits}
\def\NL{\mathop{N\!L}\nolimits}
\def\CH{\mathop{\mathrm{CH}}\nolimits}
\def\proetale{{pro\text{-}\acute{e}tale}}
\def\etale{{\acute{e}tale}}
\def\QCoh{\mathit{QCoh}}
\def\Ker{\mathop{\mathrm{Ker}}}
\def\Im{\mathop{\mathrm{Im}}}
\def\Coker{\mathop{\mathrm{Coker}}}
\def\Coim{\mathop{\mathrm{Coim}}}

% Boxtimes
%
\DeclareMathSymbol{\boxtimes}{\mathbin}{AMSa}{"02}

%
% Macros for moduli stacks/spaces
%
\def\QCohstack{\mathcal{QC}\!\mathit{oh}}
\def\Cohstack{\mathcal{C}\!\mathit{oh}}
\def\Spacesstack{\mathcal{S}\!\mathit{paces}}
\def\Quotfunctor{\mathrm{Quot}}
\def\Hilbfunctor{\mathrm{Hilb}}
\def\Curvesstack{\mathcal{C}\!\mathit{urves}}
\def\Polarizedstack{\mathcal{P}\!\mathit{olarized}}
\def\Complexesstack{\mathcal{C}\!\mathit{omplexes}}
% \Pic is the operator that assigns to X its picard group, usage \Pic(X)
% \Picardstack_{X/B} denotes the Picard stack of X over B
% \Picardfunctor_{X/B} denotes the Picard functor of X over B
\def\Pic{\mathop{\mathrm{Pic}}\nolimits}
\def\Picardstack{\mathcal{P}\!\mathit{ic}}
\def\Picardfunctor{\mathrm{Pic}}
\def\Deformationcategory{\mathcal{D}\!\mathit{ef}}
\title{An injective module whose quasi-coherent sheaf has nonvanishing H\textasciicircum{}1 on a quasi-compact open of a Noetherian affine scheme}
\author{The Stacks Project Authors}
\date{stacks-project:examples.tex:local-snapshot-20261012; GNU FDL 1.2 or later}
\begin{document}
\maketitle
\noindent\textit{Editorial note.} All delivered mathematical text is the author's own, copied from the pinned Stacks Project chapter file examples.tex; no statement or proof step was rewritten, reordered or removed.\par\medskip
\noindent\textit{Editorial note.} The counted proof is the author's complete argument: The counted proof adds only the topological-space remark Spec(A) = Spec(k[x,y]); all mathematics lives complete in the context span.\par\medskip
\noindent\textit{Editorial note.} The label lemma-nonvanishing names the counted statement, and the proof environment immediately following it in the source is delivered verbatim as its proof.\par\medskip
\noindent\textit{Editorial note.} The counted statement is carried by the author's own lemma environment in the chapter source.\par\medskip
\noindent\textit{Allowed standard prerequisites (editorial).}\begin{itemize}
\item Context span 3789-3807: with A = k[x, y, w\_1, u\_1, w\_2, u\_2, ...]/(x\textasciicircum{}n w\_n, y\textasciicircum{}n u\_n, u\_n\textasciicircum{}2, w\_n\textasciicircum{}2) and I injective, the element 1/xy of A\_\{xy\} is not in the image of I\_x  plus  I\_y -> I\_xy (clearing denominators and multiplying by u\_\{n+m\} w\_\{n+m\} gives the contradiction); U = D(x)  union  D(y) with the Mayer-Vietoris sequence 0 -> H\textasciicircum{}0(U, M\textasciitilde{}) -> M\_x  plus  M\_y -> M\_xy -> H\textasciicircum{}1(U, M\textasciitilde{}) -> 0 yields H\textasciicircum{}1(U, I\textasciitilde{}) nonzero
\item Spec(A) = Spec(k[x,y]) as topological spaces (the extra variables are square-zero artefacts)
\end{itemize}\medskip
$A_{xy} \subset I_{xy}$ is not in the image of $I_x \oplus I_y \to I_{xy}$.
Arguing by contradiction, suppose that
$$
\frac{1}{xy} = \frac{i}{x^n} + \frac{j}{y^n}
$$
for some $n \geq 1$ and $i, j \in I$. Clearing denominators we obtain
$$
(xy)^{n + m - 1} = x^my^{n + m}i + x^{n + m}y^mj
$$
for some $m \geq 0$. Multiplying with $u_{n + m}w_{n + m}$ we see
that $u_{n + m}w_{n + m}(xy)^{n + m - 1} = 0$ in $A$ which is the
desired contradiction.
Let $U = D(x) \cup D(y) \subset X = \Spec(A)$. For any $A$-module
$M$ we have an exact sequence
$$
0 \to H^0(U, \widetilde{M}) \to M_x \oplus M_y \to M_{xy} \to
H^1(U, \widetilde{M}) \to 0
$$
by Mayer-Vietoris. We conclude that $H^1(U, \widetilde{I})$ is nonzero.

\begin{lemma}
\label{lemma-nonvanishing}
There exists an affine scheme $X = \Spec(A)$ whose underlying
topological space is Noetherian and an injective
$A$-module $I$ such that $\widetilde{I}$ has nonvanishing $H^1$
on some quasi-compact open $U$ of $X$.
\end{lemma}

\begin{proof}
See above. Note that $\Spec(A) = \Spec(k[x, y])$ as topological spaces.
\end{proof}

\end{document}
